\Needspace{12\baselineskip}
\begin{center}\textit{A general (fundamental) theorem is established by examining cases, and conclusions are drawn from it.}\end{center}

\Article{130}
Having rigorously proved that every prime of the form $4n+1$, taken both positively and negatively, is a nonresidue of some smaller prime, \SourceLeaf{104}we now turn to a more exact and general comparison of primes, according to whether one is a residue or nonresidue of another.

We have rigorously proved above that $-3$ and $+5$ are residues or nonresidues of all primes that are respectively residues or nonresidues of $3, 5$.

Examining the following numbers shows that $-7$, $-11$, $+13$, $+17$, $-19$, $-23$, $+29$, $-31$, $+37$, $+41$, $-43$, $-47$, $+53$, $-59$, etc., are residues or nonresidues of all primes that, taken positively, are respectively residues or nonresidues of those primes. This examination is very easy with the help of Table II.

Anyone paying a little attention will notice that, among these primes, those of the form $4n+1$ carry a positive sign, while those of the form $4n+3$ carry a negative sign.

\Article{131}
We will soon prove that what we have found here by examining cases holds generally. Before taking up that task, however, we must work out everything that follows from the theorem if it is assumed true. We state the theorem itself as follows.

\emph{If $p$ is a prime of the form $4n+1$, then $+p$ (or, if $p$ has the form $4n+3$, then $-p$) is a residue or nonresidue of every prime that, taken positively, is a residue or nonresidue of $p$.}

Since almost everything that can be said about quadratic residues rests on this theorem, the name \emph{fundamental theorem}, which we will use from now on, is quite appropriate.

To present our reasoning as briefly as possible, we will use $a, a', a''$, etc., for primes of the form $4n+1$, and $b, b', b''$, etc., for primes of the form $4n+3$; $A, A', A''$, etc., for any numbers of the form $4n+1$, and $B, B', B''$, etc., for any numbers of the form $4n+3$. Finally, the letter $R$ between two quantities will indicate that the first is a residue of the second, while the letter $N$ will have the opposite meaning. \emph{For example,} $+5R11$, $\pm 2N5$ will mean that $+5$ is a residue of $11$, and $+2$ or $-2$ is a nonresidue of $5$. \SourceLeaf{105}Combining the fundamental theorem with the theorems of Article $111$ now gives the following propositions easily.

\[
\begin{array}{rccc}&\text{If}&&\text{then}\\[5pt]
1.&\pm aRa'&\mathinner{.\ .\ .\ .\ .\ .\ .}&\pm a' Ra\\[4pt]
2.&\pm aNa'&\mathinner{.\ .\ .\ .\ .\ .\ .}&\pm a' Na\\[4pt]
3.&\left\{\begin{array}{l}+aRb\\-aNb\end{array}\right\}&\mathinner{.\ .\ .\ .\ .\ .\ .}&\pm bRa\\[4pt]
4.&\left\{\begin{array}{l}+aNb\\-aRb\end{array}\right\}&\mathinner{.\ .\ .\ .\ .\ .\ .}&\pm bNa\\[4pt]
5.&\pm bRa&\mathinner{.\ .\ .\ .\ .\ .\ .}&\left\{\begin{array}{l}+aRb\\-aNb\end{array}\right.\\[4pt]
6.&\pm bNa&\mathinner{.\ .\ .\ .\ .\ .\ .}&\left\{\begin{array}{l}+aNb\\-aRb\end{array}\right.\\[4pt]
7.&\left\{\begin{array}{l}+bRb'\\-bNb'\end{array}\right\}&\mathinner{.\ .\ .\ .\ .\ .\ .}&\left\{\begin{array}{l}+b' Nb\\-b' Rb\end{array}\right.\\[4pt]
8.&\left\{\begin{array}{l}+bNb'\\-bRb'\end{array}\right\}&\mathinner{.\ .\ .\ .\ .\ .\ .}&\left\{\begin{array}{l}+b' Rb\\-b' Nb\end{array}\right.\\[4pt]
\end{array}
\]

\Article{132}
These contain every case arising from the comparison of two primes. The following propositions concern any numbers, but their proofs are less immediate.

\[
\begin{array}{rccc}&\text{If}&&\text{then}\\[5pt]
9.&\pm aRA&\mathinner{.\ .\ .\ .\ .\ .\ .}&\pm ARa\\[4pt]
10.&\pm bRA&\mathinner{.\ .\ .\ .\ .\ .\ .}&\left\{\begin{array}{l}+ARb\\-ANb\end{array}\right.\\[4pt]
11.&+aRB&\mathinner{.\ .\ .\ .\ .\ .\ .}&\pm BRa\\[4pt]
12.&-aRB&\mathinner{.\ .\ .\ .\ .\ .\ .}&\pm BNa\\[4pt]
13.&+bRB&\mathinner{.\ .\ .\ .\ .\ .\ .}&\left\{\begin{array}{l}-BRb\\+BNb\end{array}\right.\\[4pt]
14.&-bRB&\mathinner{.\ .\ .\ \phantom{.\ .\ .}\ .}&\left\{\begin{array}{l}+BRb\\-BNb\end{array}\right.\\[4pt]
\end{array}
\]

\SourceLeaf{106}Since the proofs of all these propositions must be drawn from the same principles, we need not develop them all. The proof of Proposition $9$, which we give, can serve as an example. First, observe that every number of the form $4n+1$ has either no factor of the form $4n+3$, or two, or four, etc.; \emph{that is,} the number of such factors (which may include equal ones) is always even. Every number of the form $4n+3$, however, contains an odd number of factors of the form $4n+3$ (\emph{that is,} either one, or three, or five, etc.). The number of factors of the form $4n+1$ remains undetermined.

\emph{Proposition} $9$ is proved as follows. Let $A$ be the product of the prime factors $a',a'',a'''$, etc., $b,b',b''$, etc. The number of factors $b,b',b''$, etc., will be even (there may also be none, which amounts to the same thing). Now, if $a$ is a residue of $A$, it is also a residue of all the factors $a',a'',a'''$, etc., $b,b',b''$, etc. So, by Propositions $1,3$ of the preceding article, each of these factors is a residue of $a$, and so is their product $A$. The same must hold for $-A$. --- If, however, $-a$ is a residue of $A$, and therefore of all the factors $a',a''$, etc., $b,b'$, etc., each of $a',a''$, etc., will be a residue of $a$, but each of $b,b'$, etc., a nonresidue. Since the number of the latter is even, the product of all the factors, \emph{that is,} $A$, will be a residue of $a$, and so will $-A$.

\Article{133}
Let us make the investigation still more general. Consider any two coprime odd numbers, with any signs, $P$ and $Q$. Think of $P$ as resolved into its prime factors without regard to its sign, and let $p$ denote how many of these have $Q$ as a nonresidue. If a prime of which $Q$ is a nonresidue occurs several times among the factors of $P$, it must also be counted several times. Similarly, let $q$ be the number of prime factors of $Q$ of which $P$ is a nonresidue. The numbers $p,q$ will then have a certain relation to each other, depending on the nature of the numbers $P,Q$. Specifically, if one of $p,q$ is even or odd, the forms of the numbers $P,Q$ will tell us whether the other is even or odd. The following table presents this relation.

$p,q$ will be both even or both odd when the numbers $P,Q$ have the forms:

\[
\begin{array}{rl}1.&+A,\quad+A'\\2.&+A,\quad-A'\\\SourceLeaf{107}3.&+A,\quad+B\\4.&+A,\quad-B\\5.&-A,\quad-A'\\6.&+B,\quad-B'\end{array}
\]

In contrast, one of the numbers $p,q$ will be even and the other odd when the numbers $P,Q$ have the forms:

\[
\begin{array}{rl}7.&-A,\quad+B\\8.&-A,\quad-B\\9.&+B,\quad+B'\\10.&-B,\quad-B'{}^{*)}\end{array}
\]

\HistoricalDisplayNote{*)}{\begin{tabular}{@{}ll@{}}Let &$l=1$ if both $P,Q\equiv3\,(\operatorname{mod.}4)$, otherwise $l=0$\\&$m=1$ if both $P,Q$ are negative, otherwise $m=0$\end{tabular}\par Then the relation depends on $l+m$.}

\emph{Example.} Let the given numbers be $-55$ and $+1197$, which belong to the fourth case. Now $1197$ is a nonresidue of one prime factor of $55$. That factor is the number $5$, while $-55$ is a nonresidue of three prime factors of $1197$, namely the numbers $3,3,19$.

If $P$ and $Q$ denote primes, these propositions reduce to those given in Article $131$. In this case $p$ and $q$ cannot be greater than $1$, so if $p$ is assumed even, it must be $=0$, \emph{that is,} $Q$ is a residue of $P$; if $p$ is odd, $Q$ is a nonresidue of $P$. And conversely. Thus, writing $a,b$ in place of $A,B$, we obtain from $8$ that if $-a$ is a residue or nonresidue of $b$, then $-b$ is a nonresidue or residue of $a$, which agrees with $3$ and $4$ of Article $131$.

In general, however, $Q$ can be a residue of $P$ only if $p=0$. So if $p$ is odd, $Q$ is certainly a nonresidue of $P$.

The propositions of the preceding article can also be derived from this without difficulty.

We will soon see, moreover, that this general presentation is more than sterile speculation, since a complete proof of the fundamental theorem could hardly be carried out without it.

\SourceLeaf{108}\Article{134}
We now turn to the deduction of these propositions.

I.\quad As before, think of $P$ as resolved into its prime factors, ignoring signs, and also resolve $Q$ into factors in any way, but taking account of the sign of $Q$. Pair each of the former factors with each of the latter. If $s$ denotes the number of all pairs in which the factor of $Q$ is a nonresidue of the factor of $P$, then $p$ and $s$ will be both even or both odd. For let the prime factors of $P$ be $f,f',f''$, etc., and among the factors into which $Q$ has been resolved, let there be $m$ that are nonresidues of $f$, $m'$ nonresidues of $f'$, $m''$ nonresidues of $f''$, etc. It is then easy to see that

\[
s=m+m'+m''+\text{etc.}
\]

while $p$ expresses how many of the numbers $m,m',m''$, etc., are odd. It follows at once that $s$ is even when $p$ is even, and odd when $p$ is odd.

II.\quad This holds generally, however $Q$ is resolved into factors. Let us turn to particular cases. First, consider the cases where one of the numbers, $P$, is positive, while the other, $Q$, has either the form $+A$ or the form $-B$. Resolve $P,Q$ into their prime factors. Give every factor of $P$ a positive sign, and every factor of $Q$ a positive or negative sign according to whether it has the form $a$ or $b$. Then $Q$ will clearly have either the form $+A$ or $-B$, as required. Pair each factor of $P$ with each factor of $Q$, and let $s$, as before, denote the number of pairs in which the factor of $Q$ is a nonresidue of the factor of $P$. Similarly, let $t$ denote the number of pairs in which the factor of $P$ is a nonresidue of the factor of $Q$. The fundamental theorem shows that these two sets of pairs are identical, so $s=t$. Finally, what we have just proved gives $p\equiv s\, (\mathrm{mod.}\ 2)$, $q\equiv t\, (\mathrm{mod.}\ 2)$, and therefore $p\equiv q\, (\mathrm{mod.}\ 2)$.

We thus have Propositions $1,3,4$ and $6$ of Article $133$.

The remaining propositions can be obtained directly by a similar method, but require one new consideration. It is easier to derive them from the preceding results as follows.

III.\quad Again, let $P,Q$ denote any two coprime odd numbers, and $p,q$ the numbers of prime factors of $P,Q$ of which $Q$, $P$, respectively, are nonresidues. Finally, let $p'$ be the number of prime factors of $P$ of which \SourceLeaf{109}$-Q$ is a nonresidue (when $Q$ itself is negative, $-Q$ clearly denotes a positive number). Divide all the prime factors of $P$ into four classes:

1) Factors of the form $a$ of which $Q$ is a residue.

2) Factors of the form $b$ of which $Q$ is a residue. Let their number be $\chi$.

3) Factors of the form $a$ of which $Q$ is a nonresidue. Let their number be $\psi$.

4) Factors of the form $b$ of which $Q$ is a nonresidue. Their number $=\omega$.

We then easily see that $p=\psi+\omega$, $p'=\chi+\psi$.

\Needspace{5\baselineskip}
Now, when $P$ has the form $\pm A$, $\chi+\omega$, and therefore also $\chi-\omega$, is even. So $p'=p+\chi-\omega\equiv p\, (\mathrm{mod.}\ 2)$. When $P$ has the form $\pm B$, a similar argument shows that the numbers $p,p'$ are incongruent modulo $2$.

IV.\quad Let us apply this to the individual cases. First, let both $P$ and $Q$ have the form $+A$. By Proposition $1$, we have $p\equiv q\, (\mathrm{mod.}\ 2)$; but $p'\equiv p\, (\mathrm{mod.}\ 2)$, so also $p'\equiv q\ (\mathrm{mod.}\ 2)$. This agrees with Proposition $2$. --- Similarly, let $P$ have the form $-A$. With $Q$ of the form $+A$, we have $p\equiv q\, (\mathrm{mod.}\ 2)$ by Proposition $2$, which we have just proved. From this, since $p'\equiv p$, we get $p'\equiv q$. Proposition $5$ is therefore proved as well.

In the same way, Proposition $7$ follows from $3$; Proposition $8$ from either $4$ or $7$; Proposition $9$ from $6$; and Proposition $10$ from the same proposition.
